\documentclass[11pt,a4paper,sans]{moderncv}
\moderncvstyle{classic}
\moderncvcolor{blue}
\usepackage[utf8]{inputenc}
\usepackage[scale=0.82]{geometry}

% moderncv compose \today en anglais par défaut ; format français sans charger babel.
\renewcommand{\today}{\number\day\space\ifcase\month\or janvier\or février\or mars\or avril\or
mai\or juin\or juillet\or août\or septembre\or octobre\or novembre\or décembre\fi\space\number\year}

\name{Alexandro}{Disla}
\title{Lettre de Motivation}
\address{Rue Saint-Louis Jeanty \#14, Route de Frère}{Port-au-Prince, Haïti}{}
\phone[mobile]{(+509) 4148-3700}
\email{alexandrodisla@hotmail.com}
\extrainfo{GitHub: \url{https://github.com/AD0791}}

\begin{document}

\recipient{Équipe de Recrutement --- Institut Panos}{Port-au-Prince, Haïti\\Poste : Responsable MEL/ACQ --- AFGHS-PANOS}
\date{\today}
\opening{Madame, Monsieur,}
\closing{Je vous prie d'agréer, Madame, Monsieur, l'expression de mes salutations distinguées,}
\enclosure[Ci-joint]{Curriculum Vitae}

\makelettertitle

Je vous adresse ma candidature au poste de Responsable Suivi, Évaluation et Apprentissage / ACQ du programme AFGHS-PANOS. Pendant trois ans, j'ai été Officier de Suivi \& Évaluation sur un programme VIH à la Caris Foundation International : j'y conduisais le rapportage de routine de cliniques multi-sites et remontais les indicateurs MER à l'USAID via DATIM, la plateforme du PEPFAR bâtie sur DHIS2. En parallèle, j'exerce comme ingénieur de données et j'administre les bases et les systèmes de collecte sur lesquels ce rapportage repose. Je réside à Port-au-Prince, le français et le créole haïtien sont mes langues maternelles, et je rédige et travaille en anglais.

Deux exigences de votre annonce sont au cœur de ce que je sais faire. La première est la réconciliation des données avec le MESI, le SISNU sur DHIS2 et iSanté Plus : c'est un problème d'ingénierie avant d'être un problème de reporting --- des définitions qui divergent d'un système à l'autre, des périodes qui ne se recouvrent pas, des patients comptés deux fois faute de résolution d'identité. J'ai construit des pipelines ETL entre sources hétérogènes et administré les bases relationnelles qui les reçoivent ; c'est mon autre métier, pas une intention. La seconde est la préparation aux vérifications indépendantes : à la Caris Foundation, j'établissais et supervisais les protocoles d'assurance qualité en remontant les chiffres rapportés jusqu'aux registres sources et en suivant les actions correctives jusqu'à clôture. Quand des jalons commandent directement les paiements du mécanisme, c'est exactement la discipline qui fait tenir un chiffre devant un vérificateur.

Sur le VIH, je veux être précis quant à mon périmètre. Mon rôle portait sur le système d'information et le cycle de rapportage d'un programme VIH --- la production, le contrôle et la soumission du jeu d'indicateurs MER, dont relèvent la rétention sous traitement et la suppression virale --- et non sur la conception clinique de ces indicateurs ni sur la PTME. Je connais le paysage national dans lequel s'inscrit ce programme, et j'ai produit des données pour DATIM sans en avoir été administrateur ; je préfère le dire clairement plutôt que de le laisser supposer.

Deux réserves, enfin, que je pose franchement. Ma scolarité en économie appliquée au C.T.P.E.A --- un cycle de quatre ans d'économie et de statistique --- est complétée, mais le Diplôme d'Études Supérieures reste en attente de la soutenance de mon mémoire de sortie : je ne revendique pas un master délivré. Et mes dix ans de suivi-évaluation et de systèmes d'information se répartissent entre cinq ans en ONG, dont trois sur le VIH, et huit ans au ministère du Plan ; je n'ai pas encore occupé de poste de conseil technique principal. Ce que j'apporte à la place est rare dans ce vivier : quelqu'un qui a rapporté dans DATIM et qui sait aussi construire le système qui l'alimente. Je reste à votre entière disposition pour un entretien.

\makeletterclosing

\end{document}
